\chapter*{Resumen}
\addcontentsline{toc}{chapter}{Resumen}

Esta tesis estudia la reconstruccion de canales EEG faltantes mediante un enfoque de procesamiento de senales sobre grafos, comparando metodos Baseline, metodos GSP instantaneos y metodos TV/tiempo bajo un protocolo experimental comun y reproducible.

El cierre B3/B4 consolido la evidencia cuantitativa con artefactos trazables de tabla final (REP-01) y significancia estadistica (STAT-02). Los resultados muestran un cruce de desempeno por metrica: Baseline presenta menor MAE y RMSE promedio, mientras TV/Tiempo presenta menor DTW promedio. Por lo tanto, la seleccion metodologica debe realizarse segun objetivo de reconstruccion y escenario de perdida, en vez de asumir una jerarquia universal.

La campana iterativa extendida it02--it130 se integra con politica GO/NO-GO: solo corridas GO se usan como evidencia confirmatoria y las NO-GO se usan para delimitar limites de validez. Adicionalmente, se mantiene la restriccion INS-13: BGSRP se reporta como proxy\_or\_partial con strict\_close=False, sin declarar equivalencia estricta 1:1 con MATLAB/GSPBox.

Como aporte principal, el trabajo entrega una base metodologica reproducible y auditable para publicacion, junto con una agenda de cierre concreta para fortalecer validacion externa y equivalencia cross-stack.
